\section{Using and Extending the Tool}
\label{sec:usage}

The same pipeline is reachable from five front ends, each aimed at one of the user groups of Section~\ref{sec:intro:audience}. None of them contains analysis logic.

\subsection{Command line}
\label{sec:usage:cli}

Installing the package (\code{pip install -e .}, Python $\ge$ 3.10) provides ten console scripts. \code{saag} runs the whole pipeline or selected steps through flags; the others run one step each:
\begin{lstlisting}
saag --all --layer system                            # every step on the imported model
saag-import   --input examples/mini_atm.json --clear # Model: import and derive DEPENDS_ON
saag-analyze  --layer mw                             # Analyze: metrics and RM scores
saag-diagnose --layer system                         # Diagnose: catalog + gate, exit 0/1/2
saag-diagnose --layer system --baseline merge_base.json --waivers waivers.yaml  # delta gate
saag-simulate fault-inject --input examples/mini_atm.json --export-json   # Simulate
saag-visualize --layer system --output output/dashboard.html             # HTML report
python cli/prescribe_graph.py --layer system         # Prescribe: verified refactoring
\end{lstlisting}
Further scripts cover generation of synthetic scenarios from YAML (\code{saag-generate}), validation (\code{saag-validate}), learned prediction (\code{saag-predict}, \code{cli/train\_graph.py}) and deployment planning (\code{saag-deploy}). The step scripts operate on the model stored in Neo4j (connection settings via flags or environment variables). Every script writes JSON, so its output can be consumed by other tools; the baseline of the delta-aware gate is simply the merge-base's \code{saag-diagnose} output, and \code{saag-diagnose}'s exit code is the CI verdict.

\subsection{Python API}
\label{sec:usage:sdk}

Researchers use the package directly. \code{saag.Pipeline} is a fluent builder whose steps run in dependency order; passing a \code{MemoryRepository} removes the need for Neo4j:
\begin{lstlisting}[language=Python]
import saag
from saag.infrastructure.memory_repo import MemoryRepository

result = (saag.Pipeline.from_json("examples/mini_atm.json", clear=True,
                                  repo=MemoryRepository())
            .analyze(layer="system")
            .simulate(layer="system", mode="exhaustive")
            .diagnose(k=5)
            .validate()
            .prescribe()
            .visualize(output="output/mini_atm.html")
            .run())
print(result.prescription.sri_improvement)
\end{lstlisting}
For step-by-step control, \code{saag.Client} exposes one method per step (\code{import\_topology}, \code{analyze}, \code{diagnose}, \code{detect\_antipatterns}, \code{simulate}, \code{validate}, \code{prescribe}, \code{deploy}, \ldots). All examples and walkthroughs in the repository use this API.

\subsection{REST service and web dashboard}
\label{sec:usage:web}

A FastAPI service (\code{api/}) exposes graph import and export, analysis, component queries, statistics, simulation, classification, validation, prediction and traffic over HTTP, with an OpenAPI description generated by FastAPI. The \code{smart/} web application (Next.js and React, charts in ECharts) is its main client. Its pages include a graph explorer that lists the model's elements and draws each one's neighbourhood with a force-directed layout coloured by type (Figure~\ref{fig:screens}a), an analysis view per layer with the RM risk profile, the criticality distribution of components and connections, a per-component score table, the detected anti-patterns and a risk topology graph coloured by criticality level (Figure~\ref{fig:screens}b), a simulator page for single and cascading failures, statistics, and validation. One Docker image starts Neo4j, the API and the web application together (\code{docker compose up}).

For settings where running a server is not an option (a CI artefact, an e-mail attachment, a review meeting), \code{saag-visualize} generates a single self-contained HTML report (\code{DashboardGenerator}) with a network view, a dependency matrix, the anti-pattern findings, the triage panel and the per-component explanations (Figure~\ref{fig:screens}c).

\begin{figure}[htbp]
\centering
\begin{tabular}{@{}c@{\hspace{3mm}}c@{}}
\screenshot{explorer.png}{0.48\textwidth} & \screenshot{analysis.png}{0.48\textwidth} \\
\small (a) Graph explorer & \small (b) Analysis view \\[2mm]
\screenshot{report.png}{0.48\textwidth} & \screenshot{prescribe_walkthrough.png}{0.48\textwidth} \\
\small (c) Self-contained HTML report & \small (d) Prescribe walkthrough
\end{tabular}
\caption{\sag{} front ends on the running example: (a, b) the \code{smart/} web application backed by Neo4j, showing the ops host's neighbourhood and the system-layer analysis (one CRITICAL and three HIGH components); (c) the static report produced by \code{saag-visualize}, here its overview tab with the per-dimension contributions to $Q(v)$; (d) one of the nine walkthrough pages (anti-patterns and prescribe), here its overview of which catalog patterns fire on the ATM scenario.}
\label{fig:screens}
\end{figure}

\subsection{Walkthroughs}
\label{sec:usage:walkthroughs}

Nine scripts in \code{examples/} each take one pipeline step apart on mini-ATM (Model, Analyze, Simulate, Predict, Diagnose, Triage, Validate, Prescribe, Deployment) and generate an interactive page from the tool's real output (Figure~\ref{fig:screens}d). They are documentation, not evidence about accuracy, but writing them was also a form of testing: several scripts recompute what they display from public helpers (the model walkthrough, for instance, recomputes every topic, component and edge weight and stops on any mismatch), and doing so surfaced several discrepancies between code and documentation (Section~\ref{sec:lessons}).

\subsection{Extension points}
\label{sec:usage:extension}

The hexagonal structure determines where extensions go:
\begin{itemize}\tightlist
  \item \emph{A new anti-pattern} is a \code{PatternSpec} entry in \code{CATALOG} (\code{saag/analysis/antipattern\_detector.py}) plus a detection method registered in the detector's dispatch table; the method receives the analysed layer and its metric statistics and returns findings. The delta gate, reports and API pick it up without change.
  \item \emph{A new storage back end} implements \code{IGraphRepository}, including the derivation rules; the simulator's derivation test then checks it against the in-memory reference.
  \item \emph{A new refactoring operator} adds a candidate type to \code{PrescriptionPolicy}, a compilation rule in \code{saag/prescription/rules.py} and a model mutation in \code{saag/prescription/mutator.py}; the verifier is operator-agnostic and applies Eq.~\eqref{eq:accept} to it unchanged.
  \item \emph{A new input source} only has to emit the JSON format of Listing~\ref{lst:json}; the real-world adapter (\code{saag/adapters/realworld\_adapter.py}) is an example.
  \item \emph{A new metric or score} is added to the structural analyzer and declared in the metric registry (\code{saag/core/metric\_registry.py}), which records for every metric whether it feeds scoring, detection, the GNN features, or only reports.
\end{itemize}
